\section{Transport through the allowed operator words}
\label{sec:word-transport}

\status{This section proves the operator-word half of incoming \texttt{MemAt}. Its input is the actual initial-symbol assertion \texttt{MemAtG}$(0,1/8)$, proved in IN9. Exact digit identities, a normal-type calculation and a numerical prefix budget prove \texttt{WordConditionAt}; linearity then gives the assertion for the whole word span. The existence of the frame and initial symbols is supplied separately.}

\subsection{The prefix budget}

Put $k=r+K$. A valid periodic frame has increasing even positions
$f_{j,0}<\cdots<f_{j,r-1}=z_j$, with
$z_{j-1}<f_{j,0}$ and $f_{j+J,i}=f_{j,i}+n'$.
Let $\mathcal W_j$ be the selected positions between $f_{j-1,0}$ and $z_j$.
The selected positions in $(z_{j-1},z_j]$, ordered as $x_{j,u}$,
are the internal letters followed by the current frame. Write $N_{j,u}$
for the first $u$ of these letters. These true letters represent the selected
residues bijectively. Define
\[
 \delta_{j,u}(s)=
 \mathbf1_{\{\exists x\in\mathcal W_j:\ x\equiv s\pmod {n'},\ x\le x_{j,u}\}}.
 \tag{WT.1}
\]
The weak inequality includes every letter affecting either adjacent
normal prefix or the plane cutoff; it also includes the previous frame
whose carry variables enter the next segment's chains. \textsf{[WT1]}

For normal joint highest weights $\kappa$, let $T_{j,u}(\kappa)$ denote
the source's \texttt{thStar}. It is the maximum of $B_{j,u}/2$ and the
active terms
$\kappa_{j,u+1,1}+m_*+s_{j,u}$ when $x_{j,u}$ is even, and
$\kappa_{j,u-1,1}+m_*+s_{j,u}$ when $u\ge1$ and $x_{j,u-1}$ is odd.
Here $s_{j,u}=\lfloor B_{j,u}/(Yk^2)\rfloor$.

Write $\mathcal M_{\alpha,\beta}(L;F)$ for the following precise universal
condition: after every actual digit/block extraction and normal joint
projection, every plane coefficient has chain orders at least
$T-\alpha s-L$ and link orders at least $T-\beta s-L$.
Positive link thresholds below these bounds are allowed; this equivalent
form permits changing one block degree during differentiation.
The input is $\mathcal M_{0,1/8}(0;F)$. \textsf{[WT1,WT6]}

\subsection{Digit extraction and the signed twist}

At one physical coordinate let $q=p^{z_j-x}$ and let $\mathcal L_\rho$
retain exponents $\rho+qb$ and replace them by $b$ in a coarse variable $y$.
For $0\le\rho<q$, direct differentiation of monomials gives
\[
\begin{aligned}
 \mathcal L_\rho(\partial_X F)
 &=\begin{cases}
 ((\rho+1)+qy\partial_y)\mathcal L_{\rho+1}F,&\rho<q-1,\\
 q\partial_y\mathcal L_0F,&\rho=q-1,
 \end{cases}\\
 \mathcal L_\rho(XF)
 &=\begin{cases}
 \mathcal L_{\rho-1}F,&\rho>0,\\
 y\mathcal L_{q-1}F,&\rho=0.
 \end{cases}
\end{aligned}
\tag{WT.2}
\]
Other residues are unchanged. After block projection, the derivative
carry term uses input degree $d+1$, and the multiplication carry term
uses $d-1$; the latter vanishes if $d=0$.
The same identities apply to the second, frame-to-carry digit extraction.
Thus the Frobenius exponents introduce scalar factors $q$, with no
amplification of the order loss. \textsf{[WT3]}

The signed-twist matrix is
$A(t)=g\operatorname{diag}(1,\ldots,t,\ldots,1)g^{-1}$.
Its entries have exponents $0,1$, and those of $A(t)^{-1}$ have exponents
$0,-1$. A vector coordinate transforms by $A(t)$ and a covector by the
inverse; the derivative transforms by the opposite matrix. Consequently
\[
\begin{array}{c|cc|cc}
 &\multicolumn{2}{c|}{\text{minimum-order loss}}
 &\multicolumn{2}{c}{\text{maximum-order increase}}\\
 &\text{vector}&\text{covector}&\text{vector}&\text{covector}\\ \hline
 y\cdot&0&1&1&0\\
 \partial_y&1&0&0&1\\
 y\partial_y&1&1&1&1
\end{array}
\tag{WT.3}
\]
Every entry is zero when the letter is outside the prefix. In particular,
Euler costs one, since one of its two factors has zero cost.
Together with (WT.2), this proves the plane derivative, multiplication
and Euler steps, including every second-digit chain. \textsf{[WT2,WT4]}

\subsection{The normal projection and one full operator}

Normal prefix projectors commute: their letter sets are nested within a
segment and disjoint across segments. On each finite multidegree space
they give a finite joint decomposition $\sum_\kappa e_\kappa=1$.
The maximum generic-twist exponent on type $\kappa_{j,u}$ is
$\kappa_{j,u,1}$. Hence for
$D\in\{\partial_y,y\cdot,y\partial_y\}$,
\[
 e_{\kappa'}D e_\kappa\ne0
 \quad\Longrightarrow\quad
 \kappa'_{j,u,1}\le
 \kappa_{j,u,1}+\mathbf1_{\{x\in N_{j,u}\}}
 \quad\text{at every cut}.                         \tag{WT.4}
\]
Indeed (WT.3) bounds the maximum before projecting, and an excluded type
has zero projection. Only these finitely many compatible input types
contribute to an output coefficient. \textsf{[WT5]}

Both adjacent indicators in $T_{j,u}$ are bounded by
$\delta_{j,u}(s)$. Taking maxima therefore gives
$T_{j,u}(\kappa')\le T_{j,u}(\kappa)+\delta_{j,u}(s)$.
For a plane operation, raise the input threshold by this same indicator
and use (WT.3). For a normal operation, use (WT.4) and the displayed
comparison. Formula (WT.2) and linearity cover all terms. Thus
\[
 \mathcal M_{0,1/8}(L;F)
 \Longrightarrow
 \mathcal M_{0,1/8}
   (L+2\delta(s);X_{s,a}\partial_{s,b}F).
 \tag{WT.5}
\]
The universal block quantifiers justify every intermediate degree change;
the complete operator preserves physical multidegree. \textsf{[WT6]}

\subsection{Absorbing the entire word}

For an allowed word $\ell$, at most $c_0$ operators occur at each selected
residue. Its accumulated loss satisfies
\[
 L_\ell(j,u)=2\sum_{(s,a,b)\in\ell}\delta_{j,u}(s)
 \le2\#\mathcal W_jc_0\le2N_1c_0.
\]
The actual frame bound gives $B_{j,u}\ge a/p^{G'}$.
The source's numerical hypothesis, for $C\ge16$, implies
$a\ge16N_1c_0Yk^2p^{G'}$; hence the integer floor satisfies
$s_{j,u}\ge16N_1c_0$.
Therefore $L_\ell\le s/8$. Induction in (WT.5), followed by weakening
the cutoffs, proves
\[
 \operatorname{MemAtG}(0,1/8;F)
 \Longrightarrow
 \operatorname{MemAtG}(1/8,3/8;\operatorname{opWord}_\ell F).
 \tag{WT.6}
\]
This is \texttt{WordConditionAt}. Each defining condition is linear in
$F$, so it extends from the initial symbols to their span and then to
the literal \texttt{OpCl} word span. Together with IN9 this proves
incoming \texttt{MemAt} for the actual family. \textsf{[WT7--WT9]}
